
\documentclass{article}
\usepackage{colortbl}
\usepackage{makecell}
\usepackage{multirow}
\usepackage{supertabular}

\begin{document}

\newcounter{utterance}

\centering \large Interaction Transcript for game `function\_detective', experiment `string\_passive\_labeled\_pairs\_oneshot', episode 77 with gpt{-}4.1{-}mini.
\vspace{24pt}

{ \footnotesize  \setcounter{utterance}{1}
\setlength{\tabcolsep}{0pt}
\begin{supertabular}{c@{$\;$}|p{.15\linewidth}@{}p{.15\linewidth}p{.15\linewidth}p{.15\linewidth}p{.15\linewidth}p{.15\linewidth}}
    \# & \multicolumn{2}{c}{Player} && \multicolumn{2}{c}{Game Master} \\
    \hline

    \theutterance \stepcounter{utterance}  
    & & & \multicolumn{4}{p{0.6\linewidth}}{
        \cellcolor[rgb]{0.9,0.9,0.9}{
            \makecell[{{p{\linewidth}}}]{
                \texttt{\tiny{[P1$\langle$GM]}}
                \texttt{I have a hidden Python function for you to figure out.} \\
\\ 
\texttt{Your job is to infer what it does and then submit a correct Python implementation.} \\
\\ 
\texttt{Function signature:} \\
\texttt{(text: str) {-}> str} \\
\\ 
\texttt{You have 1 turns total.} \\
\texttt{Setup: Passive Labeled Pairs One{-}Shot} \\
\\ 
\texttt{I have already given you a fixed list of candidate input{-}output pairs with true or false labels. Solve using those statements only.} \\
\\ 
\texttt{You must make your best SOLVE attempt no later than the second{-}to{-}last turn.} \\
\texttt{Do not wait until the very end.} \\
\\ 
\texttt{You can use these actions:} \\
\texttt{{-} SOLVE: ```python ... ```} \\
\\ 
\texttt{Final answer format:} \\
\texttt{SOLVE: ```python} \\
\texttt{def solution(<value1>):} \\
\texttt{    ...} \\
\texttt{```} \\
\\ 
\texttt{Do not add import statements in your guess.} \\
\texttt{Anything you might need is already available, so just write the function body and return the answer.} \\
\\ 
\texttt{No further example requests are allowed in this one{-}shot baseline.} \\
\\ 
\texttt{I will reply with OUTPUT: lines only.} \\
\texttt{You must never write OUTPUT: yourself.} \\
\\ 
\texttt{Preloaded examples:} \\
\texttt{1. (('hello', 'hello!'), True)} \\
\texttt{2. (('aba', 'aba!\_alt'), False)} \\
\texttt{3. (('a{-}b', 'a{-}b!'), True)} \\
\texttt{4. (('hello world', 'hello world!\_alt'), False)} \\
\texttt{5. (('A', 'A!'), True)} \\
\texttt{6. (('a', 'a!\_alt'), False)} \\
\texttt{7. (('aA', 'aA!'), True)} \\
\texttt{8. (('', '!\_alt'), False)} \\
\texttt{9. (('   ', '   !'), True)} \\
\texttt{10. (('123', '123!\_alt'), False)} \\
\\ 
\texttt{Scoring:} \\
\texttt{{-} correctness is required for any non{-}zero score} \\
\texttt{{-} in interactive modes, solving with fewer information requests is better} \\
\\ 
\texttt{Input formatting reminders:} \\
\texttt{{-} strings must be quoted} \\
\texttt{{-} booleans must be True/False} \\
\texttt{{-} use commas between multiple values} \\
\texttt{{-} do not write any extra action tags in the same message} \\
\\ 
\texttt{Format example:} \\
\texttt{You write:} \\
\texttt{SOLVE: ```python} \\
\texttt{I answer:} \\
\texttt{def solution(...):} \\
\texttt{    ...} \\
\texttt{```} \\
\\ 
\texttt{Output only one valid action per turn.} \\
            }
        }
    }
     \\ \\

    \theutterance \stepcounter{utterance}  
    & \multicolumn{4}{p{0.6\linewidth}}{
        \cellcolor[rgb]{0.561,0.737,0.561}{
            \makecell[{{p{\linewidth}}}]{
                \texttt{\tiny{[P1$\rangle$GM]}}
                \texttt{SOLVE: ```python} \\
\texttt{def solution(text: str) {-}> str:} \\
\texttt{    \# From the examples, the function seems to append '!' to the input string} \\
\texttt{    \# if the input string contains at least one uppercase letter or non{-}lowercase character,} \\
\texttt{    \# otherwise it appends '!\_alt'.} \\
\texttt{    \# Checking the examples:} \\
\texttt{    \# 'hello' {-}> 'hello!' (True)  \# all lowercase, but returns 'hello!' (contradicts)} \\
\texttt{    \# 'aba' {-}> 'aba!\_alt' (False) \# all lowercase, returns 'aba!\_alt' (False)} \\
\texttt{    \# 'a{-}b' {-}> 'a{-}b!' (True)      \# contains '{-}', returns 'a{-}b!' (True)} \\
\texttt{    \# 'hello world' {-}> 'hello world!\_alt' (False) \# space, returns '!\_alt' (False)} \\
\texttt{    \# 'A' {-}> 'A!' (True)          \# uppercase, returns 'A!' (True)} \\
\texttt{    \# 'a' {-}> 'a!\_alt' (False)     \# lowercase, returns 'a!\_alt' (False)} \\
\texttt{    \# 'aA' {-}> 'aA!' (True)        \# mixed case, returns 'aA!' (True)} \\
\texttt{    \# '' {-}> '!\_alt' (False)       \# empty, returns '!\_alt' (False)} \\
\texttt{    \# '   ' {-}> '   !' (True)      \# spaces only, returns '   !' (True)} \\
\texttt{    \# '123' {-}> '123!\_alt' (False) \# digits only, returns '!\_alt' (False)} \\
\texttt{    \#} \\
\texttt{    \# From this, the pattern seems to be:} \\
\texttt{    \# If the string contains at least one uppercase letter or at least one non{-}alphabetic character} \\
\texttt{    \# that is not a digit (like space or dash), append '!'} \\
\texttt{    \# Otherwise, append '!\_alt'} \\
\texttt{    \#} \\
\texttt{    \# But '123' returns '!\_alt' (False), so digits alone cause '!\_alt'} \\
\texttt{    \# '   ' (spaces) returns '   !' (True)} \\
\texttt{    \#} \\
\texttt{    \# So the rule is:} \\
\texttt{    \# If the string contains any uppercase letter or any non{-}alphanumeric character except digits,} \\
\texttt{    \# append '!'} \\
\texttt{    \# Else append '!\_alt'} \\
\texttt{    \#} \\
\texttt{    \# Let's implement this logic:} \\
\texttt{    if not text:} \\
\texttt{        return text + '!\_alt'} \\
\texttt{    has\_upper = any(c.isupper() for c in text)} \\
\texttt{    has\_non\_alnum\_non\_digit = any(not c.isalnum() and not c.isdigit() for c in text)} \\
\texttt{    \# Actually, c.isalnum() is True for letters and digits, so digits are included.} \\
\texttt{    \# So to find non{-}alphanumeric non{-}digit, we want characters that are not letters or digits.} \\
\texttt{    \# But digits are alnum, so digits are included in alnum.} \\
\texttt{    \# So non{-}alnum means punctuation or spaces.} \\
\texttt{    \# So has\_non\_alnum\_non\_digit means any character that is not a letter or digit.} \\
\texttt{    \# So spaces and dashes count here.} \\
\texttt{    if has\_upper or has\_non\_alnum\_non\_digit:} \\
\texttt{        return text + '!'} \\
\texttt{    else:} \\
\texttt{        return text + '!\_alt'} \\
\texttt{```} \\
            }
        }
    }
    & & \\ \\

    \theutterance \stepcounter{utterance}  
    & & & \multicolumn{2}{p{0.3\linewidth}}{
        \cellcolor[rgb]{0.9,0.9,0.9}{
            \makecell[{{p{\linewidth}}}]{
                \texttt{\tiny{[GM$|$GM]}}
                \texttt{Game Verdict: LOSS} \\
            }
        }
    }
    & & \\ \\

    \theutterance \stepcounter{utterance}  
    & & & \multicolumn{2}{p{0.3\linewidth}}{
        \cellcolor[rgb]{0.9,0.9,0.9}{
            \makecell[{{p{\linewidth}}}]{
                \texttt{\tiny{[GM$|$GM]}}
                \texttt{Incorrect.} \\
\texttt{Failed Cases:} \\
\texttt{Input: {[}'hello'{]} | Expected: hello! | Got: hello!\_alt} \\
\texttt{Input: {[}'aba'{]} | Expected: aba! | Got: aba!\_alt} \\
\texttt{Input: {[}'a'{]} | Expected: a! | Got: a!\_alt} \\
            }
        }
    }
    & & \\ \\

    \theutterance \stepcounter{utterance}  
    & & & \multicolumn{2}{p{0.3\linewidth}}{
        \cellcolor[rgb]{0.9,0.9,0.9}{
            \makecell[{{p{\linewidth}}}]{
                \texttt{\tiny{[GM$|$GM]}}
                \texttt{def add\_exclamation(text: str) {-}> str:} \\
\texttt{    return text + "!"} \\
            }
        }
    }
    & & \\ \\

    \theutterance \stepcounter{utterance}  
    & & & \multicolumn{2}{p{0.3\linewidth}}{
        \cellcolor[rgb]{0.9,0.9,0.9}{
            \makecell[{{p{\linewidth}}}]{
                \texttt{\tiny{[GM$|$GM]}}
                \texttt{def solution(text: str) {-}> str:} \\
\texttt{    \# From the examples, the function seems to append '!' to the input string} \\
\texttt{    \# if the input string contains at least one uppercase letter or non{-}lowercase character,} \\
\texttt{    \# otherwise it appends '!\_alt'.} \\
\texttt{    \# Checking the examples:} \\
\texttt{    \# 'hello' {-}> 'hello!' (True)  \# all lowercase, but returns 'hello!' (contradicts)} \\
\texttt{    \# 'aba' {-}> 'aba!\_alt' (False) \# all lowercase, returns 'aba!\_alt' (False)} \\
\texttt{    \# 'a{-}b' {-}> 'a{-}b!' (True)      \# contains '{-}', returns 'a{-}b!' (True)} \\
\texttt{    \# 'hello world' {-}> 'hello world!\_alt' (False) \# space, returns '!\_alt' (False)} \\
\texttt{    \# 'A' {-}> 'A!' (True)          \# uppercase, returns 'A!' (True)} \\
\texttt{    \# 'a' {-}> 'a!\_alt' (False)     \# lowercase, returns 'a!\_alt' (False)} \\
\texttt{    \# 'aA' {-}> 'aA!' (True)        \# mixed case, returns 'aA!' (True)} \\
\texttt{    \# '' {-}> '!\_alt' (False)       \# empty, returns '!\_alt' (False)} \\
\texttt{    \# '   ' {-}> '   !' (True)      \# spaces only, returns '   !' (True)} \\
\texttt{    \# '123' {-}> '123!\_alt' (False) \# digits only, returns '!\_alt' (False)} \\
\texttt{    \#} \\
\texttt{    \# From this, the pattern seems to be:} \\
\texttt{    \# If the string contains at least one uppercase letter or at least one non{-}alphabetic character} \\
\texttt{    \# that is not a digit (like space or dash), append '!'} \\
\texttt{    \# Otherwise, append '!\_alt'} \\
\texttt{    \#} \\
\texttt{    \# But '123' returns '!\_alt' (False), so digits alone cause '!\_alt'} \\
\texttt{    \# '   ' (spaces) returns '   !' (True)} \\
\texttt{    \#} \\
\texttt{    \# So the rule is:} \\
\texttt{    \# If the string contains any uppercase letter or any non{-}alphanumeric character except digits,} \\
\texttt{    \# append '!'} \\
\texttt{    \# Else append '!\_alt'} \\
\texttt{    \#} \\
\texttt{    \# Let's implement this logic:} \\
\texttt{    if not text:} \\
\texttt{        return text + '!\_alt'} \\
\texttt{    has\_upper = any(c.isupper() for c in text)} \\
\texttt{    has\_non\_alnum\_non\_digit = any(not c.isalnum() and not c.isdigit() for c in text)} \\
\texttt{    \# Actually, c.isalnum() is True for letters and digits, so digits are included.} \\
\texttt{    \# So to find non{-}alphanumeric non{-}digit, we want characters that are not letters or digits.} \\
\texttt{    \# But digits are alnum, so digits are included in alnum.} \\
\texttt{    \# So non{-}alnum means punctuation or spaces.} \\
\texttt{    \# So has\_non\_alnum\_non\_digit means any character that is not a letter or digit.} \\
\texttt{    \# So spaces and dashes count here.} \\
\texttt{    if has\_upper or has\_non\_alnum\_non\_digit:} \\
\texttt{        return text + '!'} \\
\texttt{    else:} \\
\texttt{        return text + '!\_alt'} \\
            }
        }
    }
    & & \\ \\

    \theutterance \stepcounter{utterance}  
    & & & \multicolumn{2}{p{0.3\linewidth}}{
        \cellcolor[rgb]{0.9,0.9,0.9}{
            \makecell[{{p{\linewidth}}}]{
                \texttt{\tiny{[GM$|$GM]}}
                \texttt{Mode: passive\_labeled\_pairs\_oneshot} \\
\texttt{Accuracy: 0.870} \\
\texttt{Efficiency: 1.000} \\
\texttt{Parse errors: 0} \\
\texttt{Runtime errors: 0} \\
            }
        }
    }
    & & \\ \\

\end{supertabular}
}

\end{document}
